\begin{ExerciseBlock}
\ExerciseHeading{07}{Auffahrunfall}{linear_kinematics_rear_end_collision}
Auf einer Straße fährt ein Auto $A$ mit konstanter Geschwindigkeit $v_0$. In einem Abstand $x_0$ vor diesem Auto beschleunigt ein Auto $B$ aus dem Stand heraus mit konstanter Beschleunigung $a > 0$.

\begin{enumerate}
	\item Stellen Sie die Ortsfunktionen $x_A(t)$ und $x_B(t)$ für die beiden Autos $A$ und $B$ auf.

	\item Zeigen Sie, dass das Auto $A$ maximal mit Geschwindigkeit \[ v_0 \leq \sqrt{2 a x_0} \] fahren darf, sodass es nicht zu einem Auffahrunfall kommt.

	\item Sollte es zum Aufprall kommen, bestimmen Sie die Aufprallgeschwindigkeit $v_\text{crash}$.

\end{enumerate}
\end{ExerciseBlock}
\UseExerciseSolution{linear_kinematics_rear_end_collision}
